% Generated draft: requires mathematical and source review.
\documentclass[11pt]{article}
\usepackage[margin=25mm]{geometry}
\usepackage{fontspec}
\setmainfont{Latin Modern Roman}
\usepackage{amsmath,amssymb,mathtools,unicode-math}
\setmathfont{Latin Modern Math}
\usepackage{enumitem,needspace,xurl,hyperref}
\setlength{\parindent}{0pt}
\setlength{\parskip}{5pt}
\setlength{\emergencystretch}{3em}
\newcommand{\N}{\mathbb{N}}
\newcommand{\Z}{\mathbb{Z}}
\newcommand{\Q}{\mathbb{Q}}
\newcommand{\R}{\mathbb{R}}
\newcommand{\C}{\mathbb{C}}
\newcommand{\F}{\mathbb{F}}
\newcommand{\A}{\mathbb{A}}
\newcommand{\PP}{\mathbb{P}}
\newcommand{\E}{\mathbb{E}}
\newcommand{\eps}{\varepsilon}
\newcommand{\dd}{\,\mathrm{d}}
\begin{document}
\section*{UnsolvedMath ID: 30000319}
\textbf{Source problem code:} catalogue entry
\subsection{Definitions and mathematical statement}
Let $G$ be the finite group, $p$ the prime, and $B$ the $p$-block intended by the source. In the usual modular setting one may take a splitting algebraically closed field $k$ of characteristic $p$ and regard $B$ as a block of $kG$. A $B$-subpair $(U,b_U)$ consists of a $p$-subgroup $U\le G$ and an appropriate block $b_U$ of $kC_G(U)$ lying under $B$ in the Brauer-pair sense. The precise coefficient system and subpair order are not given in the extraction.

Write $Z(U)$ for the center, $C_G(U)$ for the centralizer, and
\[
 N_G(U,b_U)=\{g\in N_G(U):{}^g b_U=b_U\}.
\]
For a finite group $H$, $O_p(H)$ denotes its largest normal $p$-subgroup. The restored statement singles out starting subpairs satisfying both
\[
 Z(U)\text{ is a defect group of }b_U,
 \qquad
 O_p\!\left(N_G(U,b_U)/(UC_G(U))\right)=1.
\]
These are the supplied conditions associated in the original wording with ``Alperin--Goldschmidt subpairs''; no additional condition is imported under that name.

The question is whether only chains of $B$-subpairs beginning with such a pair contribute to the alternating sums in the Ordinary Weight Conjecture. More precisely, once the report's chain indexing and summand are fixed, ask whether the full alternating sum can be restricted to those chains without changing its value. If $\mathcal C_B$ is the intended chain indexing set, $\mathcal C_B^{\mathrm{AG}}$ its specified-start subset, and $w(c)$ the intended signed summand, this means
\[
 \sum_{c\in\mathcal C_B}w(c)
 =\sum_{c\in\mathcal C_B^{\mathrm{AG}}}w(c),
\]
or cancellation of the total contribution of the complementary chains. It does not assert that each excluded chain individually has a zero summand. This notation only expresses the requested restriction: the source supplies neither actual character-count summands, signs, conjugacy conventions, chain lengths, nor any auxiliary character or defect parameters of the alternating formula.

The supplied assessment calls this a reduction/conjectural statement and discusses the projective conjecture; the exact relationship between that terminology, Dade's projective conjecture, and the Ordinary Weight Conjecture formulation must be checked in the report.
\subsection{Short English statement}
The question asks whether the Ordinary Weight Conjecture's alternating chain sums can be restricted to chains whose initial subpair satisfies the two stated local conditions. This concerns cancellation of the omitted chains as a whole, with the actual weight formula still missing.
\subsection{Sources}
\begin{itemize}
\item Local UnsolvedMath registry, permanent ID 30000319, code catalogue entry, ``Alperin--Goldschmidt Chains in Ordinary Weight Formulas.'' Proposed year 2005; proposer unspecified.
\item \emph{Cohomology of Finite Groups: Interactions and Applications} (2006), Oberwolfach Reports, pp.\ 2375--2434, DOI 10.4171/owr/2005/42; original report \url{https://doi.org/10.4171/owr/2005/42}. The displayed report year 2006 differs from the DOI and proposed year 2005.
\item Supplied reference ``A reduction theorem for Dade projective conjecture'': \url{https://ems.press/content/serial-article-files/32196}.
\item Statement status corrected\_verified; the supplied verification says the official report was checked and the wording made self-contained with normalized notation. It does not include the alternating formula itself.
\item Literature triage checked 2026-08-21: open. It says later reduction theorems do not constitute a general proof of the projective conjecture. This supplied assessment and its connection to the stated Ordinary Weight formulation require independent verification. Created 2026-08-25; updated 2026-09-16.
\item Registry curation by Ulam AI / UnsolvedMath Contributors, CC BY 4.0; underlying report and literature provenance remain as supplied. Unreviewed statement-only draft.
\end{itemize}
\end{document}
